\scanpage{8}
as recalled above\deleted{\unsure{x}}, the morphism $\underline{\operatorname{Pic}}_{Z_n/k}\handwritten{\longrightarrow}\underline{\operatorname{Pic}}_{X_o/k}$ is affine, and hence $\operatorname{NS}(Z_n)\handwritten{\longrightarrow}\operatorname{NS}(X_o)$ has finite kernel.) I expect that this is always the case, \deleted{\unsure{and it}} In case $A$ is normal of dimension $2$, this follows from Du-Val - Mumford's negative definiteness. It would he nice to have a proof in general( perhaps through a negative definiteness of an intersection matric $X^i.C^j$, where for every irreducible component $X^i$ of $X_o$, $C^i$ denotes a \inserted{suitable} curve on $X^i$, for instance defined in terms of a suitable \deleted{ample} relatively ample sheaf on $X$ by taking intersections of corresponding hypersurface sections on the $X^i$'\deleted{\unsure{s}} ?).

In case (12) is indeed injective, we have no trouble defining
\[
Q=P/\underline{D}
\tag*{(13)}
\]
as a \uline{group-scheme \unsure{locally|focally} of finite type over} $k$. We then have
\[
P^{\unsure{o|0}}\handwritten{\xrightarrow{\sim}}Q^o\ ,
\tag*{(14)}
\]
an isomorphism for the connected components, and
\[
Q/Q^o=(P/P^o)/\underline{D}\handwritten{\simeq}\underline{\operatorname{NS}}_{Z_n/k}/\underline{D}\ ,
\tag*{(15)}
\]
where $n$ is chosen large enough. Of course (14) and (11) yield an \unsure{ixomorphism|isomorphism}
\[
\underline{\operatorname{Lie}}(Q)\handwritten{\simeq}H^1(X,\underline{O}_X)\ .
\tag*{(16)}
\]
Let us remark by the way that we have an \uline{injective} restriction morphism
\[
H^1(X,\underline{O}_X)\handwritten{\hookrightarrow}H^1(U,\underline{O}_U)=H_s^2(S)
\tag*{(17)}
\]
as one sees by remarking that, with the above notations for $\underline{J}$, $\underline{L}$, \handwritten{the sheaf} \deleted{\unsure{of}} $i_{\handwritten{*}}(\underline{O}_U)/\underline{O}_X$ $(i:U\handwritten{\longrightarrow}X$ the inclusion) has a composition series whose quotients are the $\underline{L}^{-n}$, $n\handwritten{\geq}1$, whose $H^o$ is zero, and the map (17) is just the canonical map $H^1(X,\underline{O}_X)\handwritten{\unsure{\longrightarrow|\hookrightarrow}}H^1(X,i_*(\underline{O}_U))$. \unsure{Thhs|Thus}, the Lie algebra of the local Picard scheme is identified with a subspace of $H^1(U,\underline{O}_U)$. (This suggests that if $A$ is Cohen-Macaulay of dim.$\handwritten{\geq}3$, $Q$ is discrete, hence $\operatorname{Pic}(U)$ countable etc ).
% archive sheet number: 7
